\documentclass[12pt, titlepage]{article}

\usepackage{booktabs}
\usepackage{tabularx}
\usepackage{hyperref}
\hypersetup{
    colorlinks,
    citecolor=black,
    filecolor=black,
    linkcolor=red,
    urlcolor=blue
}
\usepackage[round]{natbib}

\input{../Comments}
\input{../Common}

\begin{document}

\title{Verification and Validation Report: \progname} 
\author{\authname}
\date{\today}
	
\maketitle

\pagenumbering{roman}

\section{Revision History}

\begin{tabularx}{\textwidth}{p{3cm}p{2cm}X}
\toprule {\bf Date} & {\bf Version} & {\bf Notes}\\
\midrule
Date 1 & 1.0 & Notes\\
Date 2 & 1.1 & Notes\\
\bottomrule
\end{tabularx}

~\newpage

\section{Symbols, Abbreviations and Acronyms}

\renewcommand{\arraystretch}{1.2}
\begin{tabular}{l l} 
  \toprule		
  \textbf{symbol} & \textbf{description}\\
  \midrule 
  T & Test\\
  \bottomrule
\end{tabular}\\

\wss{symbols, abbreviations or acronyms -- you can reference the SRS tables if needed}

\newpage

\tableofcontents

\listoftables %if appropriate

\listoffigures %if appropriate

\newpage

\pagenumbering{arabic}

\section{Functional Requirements Evaluation}

\subsection{AI Agent Decision Logic} 
\begin{enumerate}

\item test-FR-AI-1-1:

\textbf{Initial State:} A game state is loaded where Player A is only able to legally place a road at one specific location on the board.

\textbf{Input:} The structured GameState is sent to the AI agent via the Game State Manager.

\textbf{Expected Output:} The AI agent must select the action BuildRoad at the single remaining legal location.

\textbf{Actual Output:} The AI correctly selected the single legal road placement and executed the action without attempting any illegal moves.

\textbf{Result:} Pass

\textbf{Trace:} FR.S.3.1, FR.S.3.2, FR.S.2.1


\item test-FR-AI-1-2:

\textbf{Initial State:} GameState where Player A is 1 Victory Point away from winning and an unblocked node for the winning settlement is available.

\textbf{Input:} GameState where the immediate winning move is available.

\textbf{Expected Output:} The AI must select the action BuildSettlement at the winning node to secure the victory.

\textbf{Actual Output:} The AI selected the settlement placement at the winning node and the game concluded with Player A reaching the winning condition.

\textbf{Result:} Pass

\textbf{Trace:} FR.S.3.1, FR.S.3.2


\end{enumerate}



\subsection{User Interface and Data Management}

\begin{enumerate}



\item test-FR-UI-1-1:

\textbf{Initial State:} A completed 50-turn game with a full history stored in the accumalator.

\textbf{Input:} Request for replay including all dice rolls and development card purchases.

\textbf{Expected Output:} Successful return of UI game history with accurate logs of all actions, dice rolls, and card purchases.

\textbf{Actual Output:} System is able to retrieve and display the complete game history without errors, and the logs match the expected sequence of events.

\textbf{Result:} Pass

\textbf{Trace:} FR.S.5.1, FR.S.5.5, FR.S.7.3

\end{enumerate}

\section{Nonfunctional Requirements Evaluation}

\subsection{Performance and Scalability}

\begin{enumerate}
\item test-NFR-P-1:

\textbf{Initial State:} Full system stack deployed on the test server.

\textbf{Input/Condition:} Manually start 1, 2, and 3 concurrent games with multiple players, performing typical actions (build, trade, roll dice) in each session.

\textbf{Expected Output:} The system remains responsive; actions in each game update within acceptable time (no noticeable lag), and no crashes or errors occur.

\textbf{Actual Output:} System remained responsive in all three concurrent sessions; UI and AI bot updates were consistent and timely.

\textbf{Result:} Pass

\textbf{Trace:} NFR.S.1
\end{enumerate}

\subsection{Usability}

\begin{enumerate}
\item test-NFR-U-1:

\textbf{Initial State:} Final prototype deployed in a browser.

\textbf{Input/Condition:} 5--7 participants familiar with Catan play a full game against the AI bot and explore all interface options (build, trade, end turn).

\textbf{Expected Output:} Players can understand the interface and perform actions without guidance; overall experience is intuitive.

\textbf{Actual Output:} All participants successfully completed games. Users rated clarity, navigation, and understanding of AI bot behavior as 4.3/5 on average. Some suggestions for improving tooltip visibility were noted.

\textbf{Result:} Pass

\textbf{Trace:} NFR.S.2

\item test-NFR-U-2:

\textbf{Initial State:} UI opened in Chrome/Firefox with accessibility tools enabled.

\textbf{Input/Condition:} Apply simulated color-blindness filters and vary browser zoom levels between 80\% and 150\%.

\textbf{Expected Output:} All critical UI elements remain legible; actions can still be performed correctly.

\textbf{Actual Output:} Text, buttons, and board elements remained readable and interactive under all tested conditions.

\textbf{Result:} Pass

\textbf{Trace:} NFR.S.2
\end{enumerate}

\subsection{Maintainability}

\begin{enumerate}
\item test-NFR-M-1:

\textbf{Initial State:} Complete codebase committed to repository.

\textbf{Input/Condition:} Run static analysis (flake8, SonarQube) and perform peer review of module structure (AI bot, CV, UI, Database).

\textbf{Expected Output:} No high-severity issues; code is modular and well-documented for easy updates.

\textbf{Actual Output:} Static analysis reported only minor style warnings. Peer review confirmed clear module separation and sufficient documentation.

\textbf{Result:} Pass

\textbf{Trace:} NFR.S.3
\end{enumerate}

\subsection{Installability}

\begin{enumerate}
\item test-NFR-I-1:

\textbf{Initial State:} Packaged release candidate available.

\textbf{Input/Condition:} Install and launch the system on Windows 11 and macOS using Chrome and Firefox.

\textbf{Expected Output:} Application installs and runs correctly on all platforms.

\textbf{Actual Output:} Successful installation and execution on all tested platforms. Minor permission prompts appeared but did not affect functionality.

\textbf{Result:} Pass

\textbf{Trace:} NFR.S.4
\end{enumerate}

\subsection{Data Integrity}

\begin{enumerate}
\item test-NFR-D-1:

\textbf{Initial State:} Active game session connected to the database.

\textbf{Input/Condition:} Perform multiple manual game actions and simulate a sudden network interruption during one action.

\textbf{Expected Output:} The database remains consistent; interrupted updates do not corrupt data.

\textbf{Actual Output:} Interrupted update rolled back successfully. All game states remained valid and accurate.

\textbf{Result:} Pass

\textbf{Trace:} NFR.S.5

\item test-NFR-D-2:

\textbf{Initial State:} Active game session.

\textbf{Input/Condition:} Attempt to input invalid game actions manually (e.g., building where not allowed, invalid resource quantities).

\textbf{Expected Output:} Invalid actions are rejected by the system; user is informed of the error.

\textbf{Actual Output:} All invalid actions were rejected; logs correctly recorded failures.

\textbf{Result:} Pass

\textbf{Trace:} NFR.S.5
\end{enumerate}

\subsection{Availability}

\begin{enumerate}
\item test-NFR-AV-1:

\textbf{Initial State:} System running with an active game.

\textbf{Input/Condition:} Manually terminate the server process mid-game and restart.

\textbf{Expected Output:} System restores previous game state within one minute, allowing the game to continue without data loss.

\textbf{Actual Output:} Service restarted in 32 seconds; last committed game state restored and session resumed successfully.

\textbf{Result:} Pass

\textbf{Trace:} NFR.S.6
\end{enumerate}

	
\section{Comparison to Existing Implementation}	

This section will not be appropriate for every project.

\section{Unit Testing}

This section outlines the unit testing strategy and details the specific test cases used to verify individual components of the RLCatan system.
The automated unit tests focus on ensuring that API endpoints, internal path resolution logic, and analysis interfaces function correctly in isolation.

\subsection{Unit Testing Scope}

The scope of these unit tests covers the backend API endpoints, path resolution utilities, and the Monte Carlo Tree Search (MCTS) analysis integration.
Modules and underlying models developed by external dependencies are out of scope for these unit tests and are assumed to be verified by their respective maintainers.

\subsection{Tests for Functional Requirements}

\subsubsection{API and Game Management Module}

\begin{enumerate}

\item{\texttt{test_get_bots_endpoint}\}
Type: Functional, Dynamic, Automatic

Initial State: The API server is initialized and ready to receive requests.

Input: A simple HTTP GET request to the \texttt{/api/bots} URL.

Output: A 200 OK status code and a JSON response containing a list of available bot objects.

Test Case Derivation: Ensures the bot registry is populated correctly by checking that the returned list includes standard bot IDs like \texttt{catanatron_ab_2}, \texttt{random}, and \texttt{human}, and confirms that the total number of available bots is at least 10.

How test will be performed: Automated unit test using an HTTP test client.

\item{\texttt{test_stress_test_endpoint}\}
Type: Functional, Dynamic, Automatic

Initial State: The API server is initialized.

Input: An HTTP GET request to the \texttt{/api/stress-test} endpoint.

Output: A 200 OK status code and a JSON object representing the game state.

Test Case Derivation: Verifies that the system can successfully create and advance a game instance without crashing by validating that the response contains essential game data structures, specifically "nodes" and "edges".

How test will be performed: Automated unit test using an HTTP test client.

\item{\texttt{test_player_factory_custom_bot}\}
Type: Functional, Dynamic, Automatic

Initial State: Internal functions are mocked to simulate a custom bot named "my_custom_bot".

Input: A POST request to \texttt{/api/games} with "BOT:my_custom_bot" in the player list.

Output: A successful 200 OK response containing a \texttt{game_id}.

Test Case Derivation: Confirms that the backend correctly identifies the custom bot prefix, retrieves the mocked configuration, and attempts to instantiate the appropriate player class (like PPOPlayer) with the correct model path.

How test will be performed: Automated unit test with mocked dependencies.

\item{\texttt{test_create_game_all_player_types}\}
Type: Functional, Dynamic, Automatic

Initial State: The API server is initialized.

Input: A series of POST requests to \texttt{/api/games}, where each request uses a different player type (e.g., CATANATRON, VALUE_FUNCTION, HUMAN) in the configuration.

Output: A 200 OK status code and a valid \texttt{game_id} for each request.

Test Case Derivation: Serves as a comprehensive smoke test to verify that every registered player class is correctly implemented and can be initialized by the game engine.

How test will be performed: Automated unit test iterating through all supported built-in player types.

\end{enumerate}

\subsubsection{Path Resolution Module}

This module tests the utility functions responsible for safely and accurately resolving file paths for model assets across different deployment environments.

\begin{enumerate}

\item{\texttt{test_resolve_model_path_absolute_hpc}\}
Type: Functional, Dynamic, Automatic

Initial State: Path resolution utilities are loaded.

Input: A specific string representing an absolute file path typical of an HPC setup (e.g., \texttt{/nfs/features/rlcatan/...}).

Output: A normalized string path that has been rebased to the application's local directory structure (e.g., starting with \texttt{/app}).

Test Case Derivation: Ensures that absolute paths from training environments are correctly mapped to valid paths in the deployment environment.

How test will be performed: Automated unit test.

\item{\texttt{test_resolve_model_path_relative}\}
Type: Functional, Dynamic, Automatic

Initial State: Path resolution utilities are loaded.

Input: A standard relative path string (e.g., \texttt{data/models/v2}).

Output: An absolute path that correctly appends the input to the application's root directory.

Test Case Derivation: Confirms the system can locate local model files correctly by asserting that the resolved path ends with the expected \texttt{app/data/models/v2} suffix.

How test will be performed: Automated unit test.

\end{enumerate}

\subsubsection{Game Analysis Module}

This module verifies the integration with the Monte Carlo Tree Search (MCTS) engine to ensure active games can be properly analyzed.

\begin{enumerate}

\item{\texttt{test_mcts_analysis_endpoint}\}
Type: Functional, Dynamic, Automatic

Initial State: A valid game is created via the API, and the GameAnalyzer is mocked to return fixed win probabilities.

Input: A GET request to that game's analysis URL (\texttt{.../mcts-analysis}).

Output: A 200 OK status and a JSON response where success is True.

Test Case Derivation: Proves the endpoint correctly interfaces with the analysis engine by validating that the returned win probabilities match the mocked values (e.g., RED=0.5).

How test will be performed: Automated unit test with a mocked GameAnalyzer.

\item{\texttt{test_mcts_analysis_endpoint_not_found}\}
Type: Functional, Dynamic, Automatic

Initial State: The API server is initialized.

Input: A GET request to the analysis endpoint with a non-existent game ID (e.g., 99999).

Output: An error status code, specifically checking for either a 404 Not Found or a 500 Internal Server Error.

Test Case Derivation: Confirms that the application handles invalid game IDs without crashing or returning a false success.

How test will be performed: Automated unit test using an HTTP test client.

\end{enumerate}

\section{Changes Due to Testing}

\wss{This section should highlight how feedback from the users and from 
the supervisor (when one exists) shaped the final product.  In particular 
the feedback from the Rev 0 demo to the supervisor (or to potential users) 
should be highlighted.}

\section{Automated Testing}
		
\section{Trace to Requirements}
		
\section{Trace to Modules}		

\section{Code Coverage Metrics}

\bibliographystyle{plainnat}
\bibliography{../../refs/References}

\newpage{}
\section*{Appendix --- Reflection}

The information in this section will be used to evaluate the team members on the
graduate attribute of Reflection.

\input{../Reflection.tex}

\begin{enumerate}
  \item What went well while writing this deliverable? 
  \item What pain points did you experience during this deliverable, and how
    did you resolve them?
  \item Which parts of this document stemmed from speaking to your client(s) or
  a proxy (e.g. your peers)? Which ones were not, and why?
  \item In what ways was the Verification and Validation (VnV) Plan different
  from the activities that were actually conducted for VnV?  If there were
  differences, what changes required the modification in the plan?  Why did
  these changes occur?  Would you be able to anticipate these changes in future
  projects?  If there weren't any differences, how was your team able to clearly
  predict a feasible amount of effort and the right tasks needed to build the
  evidence that demonstrates the required quality?  (It is expected that most
  teams will have had to deviate from their original VnV Plan.)
\end{enumerate}

\end{document}